\chapter{Đánh giá sai số, chuỗi và FPGA}
\label{ch:measurements}
\section{Sai số số học và thời gian thực thi}

\subsection{Sai số của riêng bộ nhân lặp}
Để đánh giá riêng khối nhân trước khi đóng gói và cộng tích lũy, hệ thống kiểm
chứng đọc đồng thời giá trị nguyên $M$, dấu, số mũ cơ sở, số vòng thực thi và
hai cờ cắt tại thời điểm \texttt{core\_done} (trước chuẩn hóa và \texttt{Packer});
với trường hợp $C=0$, từ mã và cờ đầu ra tiếp tục được đối chiếu hoàn toàn với
mô hình chuẩn. Hệ thống \texttt{ModelSim} đã xác nhận không phát sinh sai lệch trên
toàn bộ $1.634.048$ giao dịch trải đều qua $13$ cấu hình giới hạn vòng $n \in [0, 12]$.
Bộ mẫu cố định gồm $125.696$ cặp cho mỗi mức $n$ bao gồm $100.000$ cặp lấy mẫu
đều hữu hạn khác $0$, $25.600$ cặp phân tầng theo \textit{regime}/dấu, $68$ cặp
chủ đích biên, $13$ cặp \texttt{Zero} và $15$ cặp \texttt{NaR}; phép thống kê
sai số tương đối chỉ tính trên các tích hữu hạn khác $0$.

Gọi $P = AB$ là tích chính xác giữ đủ $26$ bit fraction đầu vào, $P_{12}$ là
tích chính xác sau khi cắt fraction về $12$ bit và $\widehat P$ là giá trị tích
từ RTL. Ba thước đo dưới đây cho phép tách bạch ảnh hưởng giữa việc cắt ngõ vào
và việc giới hạn số vòng lặp:
\[
E_{\text{tổng}}=100\frac{|\widehat P-P|}{|P|},\qquad
E_{\text{đầu vào}}=100\frac{|P_{12}-P|}{|P|},\qquad
E_{\text{vòng}}=100\frac{|\widehat P-P_{12}|}{|P_{12}|}.
\]
Phân vị $\mathrm{P95}$ được xác định tại vị trí $\lceil 0{,}95N\rceil$ sau khi
sắp xếp tăng dần. Ở cấu hình $n=12$, tích RTL trùng khớp hoàn toàn với mốc $P_{12}$
trên cả $125.668$ cặp hữu hạn khác $0$ ($\widehat P = P_{12}$). Trên $100.000$
cặp lấy mẫu đều, $\mathrm{P95}$ và cực đại của sai số cắt đầu vào $E_{\text{đầu vào}}$
cố định ở mức \textbf{0,0297\%} và \textbf{0,0460\%} trên toàn bộ $13$ mức $n$.

\begin{table}[htbp]
  \centering
  \footnotesize
  \setlength{\tabcolsep}{4pt}
  \caption{Sai số trước Packer và chu kỳ thực thi trên 100.000 cặp $A/B$ lấy mẫu đều ở cấu hình \texttt{OPS1} (3 chữ số có nghĩa cho sai số).}
  \label{tab:prepack}
  \begin{tabular}{rrrrrrrr}
    \toprule
    $n$ & $E_{\text{tổng}}$ P95 & $E_{\text{tổng}}$ max &
    $E_{\text{vòng}}$ P95 & $E_{\text{vòng}}$ max & Vòng TB & Vòng max & Chu kỳ F TB \\
    \midrule
    0  & 47,1\% & 50,0\% & 47,1\% & 50,0\% & 0,00 & 0  & 14,00 \\
    1  & 20,7\% & 25,0\% & 20,7\% & 25,0\% & 1,00 & 1  & 14,00 \\
    2  & 7,60\% & 12,5\% & 7,59\% & 12,5\% & 1,99 & 2  & 14,99 \\
    3  & 2,54\% & 6,21\% & 2,53\% & 6,20\% & 2,95 & 3  & 15,95 \\
    4  & 0,759\% & 3,09\% & 0,746\% & 3,08\% & 3,81 & 4  & 16,81 \\
    5  & 0,212\% & 1,52\% & 0,198\% & 1,51\% & 4,46 & 5  & 17,46 \\
    6  & 0,0598\% & 0,742\% & 0,0429\% & 0,733\% & 4,84 & 6  & 17,84 \\
    7  & 0,0320\% & 0,351\% & 0 & 0,342\% & 4,98 & 7  & 17,98 \\
    8  & 0,0299\% & 0,156\% & 0 & 0,147\% & 5,02 & 8  & 18,02 \\
    9  & 0,0297\% & 0,0582\% & 0 & 0,0489\% & 5,03 & 9  & 18,03 \\
    10 & 0,0297\% & 0,0460\% & 0 & 0 & 5,03 & 10 & 18,03 \\
    11 & 0,0297\% & 0,0460\% & 0 & 0 & 5,03 & 10 & 18,03 \\
    12 & 0,0297\% & 0,0460\% & 0 & 0 & 5,03 & 10 & 18,03 \\
    \bottomrule
  \end{tabular}
  \par\smallskip
  \footnotesize Cột \textit{Chu kỳ F TB} tính từ công thức $13+\max(1,t)$ (kể cả $38$ ca có $t=0$ ở mẫu lấy đều). Ca chủ đích ở $n=12$ đạt tối đa $12$ vòng ($25$ chu kỳ) khi cả hai phần định trị toàn bit~1.
\end{table}

Bảng~\ref{tab:prepack} trình bày diễn biến sai số và số chu kỳ trung bình của
\texttt{MAC-F} khi tăng giới hạn vòng lặp $n$ trên $100.000$ cặp lấy mẫu đều.
Nhờ cơ chế \texttt{OPS1} kết hợp dừng sớm, từ $n=8$ lên $n=12$, số chu kỳ trung
bình của \texttt{MAC-F} gần như giữ nguyên ($18{,}02$ so với $18{,}03$ chu kỳ)
và $\mathrm{P95}$ đã bão hòa sát ngưỡng cắt 12 bit ($0{,}0299\%$ xuống $0{,}0297\%$),
nhưng việc mở rộng lên $n=12$ giúp cắt giảm \textbf{3,39 lần} sai số cực đại
quan sát (từ $0{,}156\%$ xuống $0{,}0460\%$) bằng cách loại bỏ toàn bộ các mẫu
còn sót số hạng. Vì vậy đồ án chọn \texttt{OPS1/n12} làm cấu hình mặc định ưu tiên
độ chính xác, trong khi $n=6\text{--}8$ phù hợp làm chế độ tiết kiệm thời gian
chặn biên trên khi gặp dữ liệu toàn bit~1 (cắt trần độ trễ từ $25$ xuống
$19\text{--}21$ chu kỳ). Kết quả chi tiết cho nhóm phân tầng và các ca biên chủ
địch được lưu riêng trong bộ dữ liệu kiểm chứng đi kèm.
\input{figures/prepack_error_curve}

Để định vị phương pháp đánh giá trong bối cảnh nghiên cứu trước đây, cần lưu ý rằng
Norris và Kim~\cite{norris2021iterativeposit} dùng chuẩn SPFP chuyển sang Posit
với độ dài run regime $m \le 15$ cho Bảng~I và DPFP cho miền $m \ge 18$ ở Bảng~II,
đối chiếu kết quả đầu ra đã đóng gói trên $200$ triệu phép nhân. Trong đồ án này,
phép đo trước \texttt{Packer} kết hợp cùng hai mốc $P_{12}$ và $P$ cho phép tách
bạch trực tiếp sai số thuật toán lặp khỏi sai số lượng tử hóa đầu ra.

\subsection{Sai số toàn MAC khi ghép cặp NF và F}
Thử nghiệm ghép cặp tái sử dụng tập $A/B$ trên bằng cách cấp thêm ba kịch bản toán
hạng cộng gồm $C = R(P/2)$, $R(-P/2)$ và $R(-P)$ đối với các cặp hữu hạn khác $0$
(cấp $C=0$ khi $A$ hoặc $B$ là \texttt{Zero}/\texttt{NaR}). Cả hai lõi \texttt{MAC-NF}
và \texttt{MAC-F} cùng vận hành tại cấu hình \texttt{OPS1/n12} trên $377.032$ bộ
$(A,B,C)$ và khớp $100\%$ với mô hình vàng của từng kiến trúc. Sai số toàn MAC được
đối chiếu trực tiếp với tổng lý tưởng đủ bit $r = P + C$ và phân tách riêng hai vùng
triệt tiêu gồm $r = 0$ và $0 < |r| \le 2^{-12}|P|$.

\begin{table}[htbp]
  \centering
  \small
  \caption{Sai số tương đối (\%) trên nhóm $100.000$ cặp $A/B$ lấy mẫu đều ở \texttt{OPS1/n12} (3 chữ số có nghĩa).}
  \label{tab:paired-mac}
  \begin{tabular}{lrrrr}
    \toprule
    C cấp vào & P95 NF & P95 F & Max NF & Max F \\
    \midrule
    $R(P/2)$  & 0,0198 & 0,0198 & 0,256 & 0,256 \\
    $R(-P/2)$ & 0,0595 & 0,0594 & 0,765 & 0,576 \\
    \bottomrule
  \end{tabular}
\end{table}

Trên các trường hợp $C = R(P/2)$ và $R(-P/2)$ (Bảng~\ref{tab:paired-mac}), phân vị
$\mathrm{P95}$ giữa \texttt{MAC-NF} và \texttt{MAC-F} gần như tương đương, song
lõi hợp nhất \texttt{MAC-F} giảm rõ rệt sai số cực đại ở nhánh trái dấu từ $0{,}765\%$
xuống $0{,}576\%$ do tránh làm tròn trung gian trước khi trừ. Ở kịch bản ép triệt tiêu
$C = R(-P)$ gồm $125.668$ mẫu hữu hạn, phép đo sử dụng hai chỉ số bất biến theo
thang là khoảng cách mã (\texttt{code-distance}) và sai số tương đối chuẩn hóa theo
tích $|err|/|P|$ để tránh hiện tượng thổi phồng giá trị tuyệt đối trên dải động rộng
(Bảng~\ref{tab:cancel-target-metrics}). Kết quả tách lớp cho thấy $4.272$ mẫu có tổng
lý tưởng $r = 0$ đều thuộc nhóm \texttt{sum\_class=zero}; cả hai lõi cùng trả đúng
\texttt{0x00000000} ở $3.966$ ca và hòa nhau ở toàn bộ $4.272$ ca ($306$ ca lệch
$0$ giống nhau bắt nguồn từ bước cắt 12 bit đầu vào của bộ nhân). Trong $108.703$
ca gần $0$ ($r \ne 0$), trung vị $|err|/|P|$ của NF và F lần lượt đạt $1{,}63{\times}10^{-4}$
và $1{,}62{\times}10^{-4}$, nhất quán với ngưỡng xấp xỉ $12$ bit. Ngược lại, trên
$12.693$ mẫu thuộc lớp \texttt{regular} (có $2.881$ ca kẹp cực tiểu và $3.328$ ca
kẹp cực đại ở $R(-P)$), \texttt{MAC-NF} chịu sai lệch kẹp ngưỡng khiến trung vị
\texttt{code-distance} vọt lên $2{,}15{\times}10^9$ và đỉnh $|err|/|P|$ chạm $2^{240}$
($1{,}77{\times}10^{72}$), trong khi \texttt{MAC-F} giữ trung vị khoảng cách mã bằng
$0$ và đỉnh $|err|/|P|$ ở mức $1{,}00$ (vượt trội trên $12.691/12.693$ mẫu theo
\texttt{code-distance}).

\begin{table}[!htbp]
  \centering
  \caption{Tóm tắt sai số nhóm $C=R(-P)$ ($n=125.668$, 3 chữ số có nghĩa; bảng đầy đủ tại Phụ lục~B).}
  \label{tab:cancel-target-metrics}
  \footnotesize
  \setlength{\tabcolsep}{4pt}
  \renewcommand{\arraystretch}{1.12}
  \begin{tabularx}{\textwidth}{@{}>{\raggedright\arraybackslash}p{30mm} r *{2}{>{\raggedleft\arraybackslash}X} >{\raggedleft\arraybackslash}p{26mm}@{}}
    \toprule
    \textbf{Tập con ($C=R(-P)$)} & \textbf{Số mẫu} & \textbf{P50 NF / F} & \textbf{P95 NF / F} & \textbf{F thắng / hòa / thua} \\
    \midrule
    \multicolumn{5}{@{}l}{\textbf{Khoảng cách mã} $\texttt{code-distance} = |\mathrm{code}_s(d) - \mathrm{code}_s(R(r))|$ (đơn vị bước mã)} \\
    \midrule
    $r = 0$ (tách riêng) & 4.272 & $0$ / $0$ & $2{,}56{\times}10^7$ / $2{,}56{\times}10^7$ & 0 / 4.272 / 0 \\
    \texttt{near\_zero} ($r \ne 0$) & 108.703 & $2{,}99{\times}10^8$ / $2{,}94{\times}10^8$ & $3{,}31{\times}10^9$ / $3{,}37{\times}10^9$ & 13.053 / 83.564 / 12.086 \\
    \texttt{regular} (có \texttt{clamp}) & 12.693 & $2{,}15{\times}10^9$ / $0$ & $2{,}15{\times}10^9$ / $75$ & 12.691 / 0 / 2 \\
    \textbf{Toàn nhóm} $C=R(-P)$ & 125.668 & $2{,}81{\times}10^8$ / $2{,}06{\times}10^8$ & $3{,}08{\times}10^9$ / $3{,}13{\times}10^9$ & 25.744 / 87.836 / 12.088 \\
    \midrule
    \multicolumn{5}{@{}l}{\textbf{Sai số tương đối theo độ lớn tích} $|err|/|P| = |d - (P+C)| / |P|$} \\
    \midrule
    $r = 0$ (tách riêng) & 4.272 & $0$ / $0$ & $7{,}09{\times}10^{-5}$ / $7{,}09{\times}10^{-5}$ & 0 / 4.272 / 0 \\
    \texttt{near\_zero} ($r \ne 0$) & 108.703 & $1{,}63{\times}10^{-4}$ / $1{,}62{\times}10^{-4}$ & $2{,}96{\times}10^{-4}$ / $2{,}96{\times}10^{-4}$ & 12.753 / 83.578 / 12.372 \\
    \texttt{regular} (có \texttt{clamp}) & 12.693 & $0{,}932$ / $0{,}991$ & $1{,}33{\times}10^{36}$ / $1{,}00$ & 12.473 / 0 / 220 \\
    \textbf{Toàn nhóm} $C=R(-P)$ & 125.668 & $1{,}70{\times}10^{-4}$ / $1{,}64{\times}10^{-4}$ & $0{,}981$ / $0{,}998$ & 25.226 / 87.850 / 12.592 \\
    \bottomrule
  \end{tabularx}
  \vspace{2pt}
  {\footnotesize Ghi chú: Toàn bộ tập $r=0$ gồm 4.272 mẫu đều thuộc \texttt{sum\_class=zero}, trong đó 3.966 mẫu khớp chính xác $0$ và 306 mẫu có sai khác giống nhau trên cả NF và F do cắt bỏ bit thấp của $P$ ở bộ nhân 12 bit trước khi cộng $C$.}
\end{table}


Tính chung toàn bộ $377.017$ phép thử hữu hạn ghép cặp, \texttt{MAC-F} cho sai số
nhỏ hơn \texttt{MAC-NF} ở $47.534$ trường hợp, bằng nhau ở $305.545$ trường hợp và
lớn hơn ở $23.938$ trường hợp. Hiện tượng $23.938$ ca ($6{,}3\%$) mà \texttt{MAC-NF}
có sai số tương đối đơn bước nhỏ hơn bắt nguồn từ \textbf{thiên lệch cắt một phía}
của bộ nhân lặp: như đã kiểm chứng ở Chương~\ref{ch:arithmetic}, tích $\widehat{P}$
luôn nhỏ hơn hoặc bằng tích lý tưởng $P$ trên $100\%$ mẫu hữu hạn ($125.668/125.668$
cặp tại \texttt{OPS1/n12}). Do đó, bước làm tròn RNE trung gian $R(\widehat{P})$ ở
\texttt{MAC-NF} khi làm tròn lên sẽ vô tình triệt tiêu bớt một phần thiên lệch âm
đó trong từng phép đơn lẻ (như tại $A=\texttt{1141ee15}$, $B=\texttt{33f45b16}$,
$C=\texttt{0c9e9823}$ nơi cả hai cùng đạt $\approx 0{,}00786\%$ khi làm tròn 3 chữ số
có nghĩa, còn độ lệch gốc của \texttt{MAC-NF} hơi thấp hơn). Dù vậy, xu hướng toàn cục
vẫn nghiêng mạnh về \texttt{MAC-F} (số ca thắng gấp đôi số ca thua), đồng thời
\texttt{MAC-F} tránh được sai lệch kẹp ngưỡng cực đại tới $2^{240}$ trên nhóm
\texttt{regular} của $C=R(-P)$ (dù vẫn ghi nhận trung thực $2$ ca trên \texttt{code-distance}
và $220$ ca trên $|err|/|P|$ mà lượng tử hóa trung gian của \texttt{MAC-NF} bù ngẫu
nhiên tốt hơn). Cần lưu ý rằng lớp \texttt{regular} chỉ xuất hiện khi $|P|$ vượt
quá dải biểu diễn của \texttt{Posit(32,3)} nên phép lấy $R(-P)$ bị kẹp vào giá
trị cực đại/cực tiểu; đây là ca biên do cách dựng tập thử tổng hợp. Trên dữ liệu
biên độ thông thường (như $C=R(P/2)$ và các bài toán chuẩn ở Mục~\ref{sec:app-benchmarks}),
\texttt{MAC-NF} và \texttt{MAC-F} đạt độ chính xác gần như tương đương.

\subsection{Chuỗi tích lũy và thời gian thực thi theo xung nhịp thực}
Phép đánh giá chuỗi sử dụng $48$ chuỗi $A/B$ cố định chia đều cho ba tính chất
(cùng dấu, xen dấu và gần triệt tiêu) ở bốn độ dài $L \in \{16, 64, 256, 1.024\}$
tại cấu hình \texttt{OPS1/n12} với điều kiện ngõ ra luôn sẵn sàng. Hai lõi
\texttt{MAC-NF} và \texttt{MAC-F} khởi tạo $C=0$ rồi hồi tiếp kết quả liên tiếp,
trong khi \texttt{MAC-Quire} khởi tạo bằng \texttt{CLEAR}, cộng tích lên thanh ghi
972 bit và đọc kết quả cuối bằng lệnh \texttt{READ}. Số chu kỳ đo trực tiếp trên
\texttt{ModelSim} qua $16.320$ phép MAC mỗi lõi (cộng $96$ lệnh điều khiển quire)
khớp $100\%$ với mô hình tham chiếu.

\begin{table}[!htbp]
  \centering
  \footnotesize
  \setlength{\tabcolsep}{3.5pt}
  \caption{Trung vị số chu kỳ (của 4 chuỗi/nhóm) và thời gian thực thi quy đổi theo tần số cực đại sau định tuyến.}
  \label{tab:chain-cycles}
  \begin{tabular}{lrrrrrrr}
    \toprule
    Nhóm chuỗi & $L$ &
    \multicolumn{3}{c}{Số chu kỳ (trung vị)} &
    \multicolumn{3}{c}{Thời gian thực tại $F_{\max,\mathrm{R2R}}$ ($\mu$s)} \\
    \cmidrule(lr){3-5}\cmidrule(lr){6-8}
    & & NF & F & Quire & \begin{tabular}[t]{@{}r@{}}NF\\(135,9\,MHz)\end{tabular} & \begin{tabular}[t]{@{}r@{}}F\\(154,4\,MHz)\end{tabular} & \begin{tabular}[t]{@{}r@{}}Quire\\(98,2\,MHz)\end{tabular} \\
    \midrule
    Cùng dấu & 16 & 345,5 & 297,5 & 727,5 & 2,54 & 1,93 & 7,41 \\
    Cùng dấu & 64 & 1.348,5 & 1.156,5 & 2.690,5 & 9,93 & 7,49 & 27,4 \\
    Cùng dấu & 256 & 5.389,5 & 4.621,5 & 10.571,5 & 39,7 & 29,9 & 108 \\
    Cùng dấu & 1.024 & 21.528 & 18.456 & 42.070 & 158 & 120 & 429 \\
    Xen dấu & 16 & 335 & 287 & 717 & 2,47 & 1,86 & 7,30 \\
    Xen dấu & 64 & 1.328 & 1.136 & 2.670 & 9,78 & 7,36 & 27,2 \\
    Xen dấu & 256 & 5.370,5 & 4.602,5 & 10.552,5 & 39,5 & 29,8 & 108 \\
    Xen dấu & 1.024 & 21.552 & 18.480 & 42.094 & 159 & 120 & 429 \\
    Gần triệt tiêu & 16 & 335 & 287 & 717 & 2,47 & 1,86 & 7,30 \\
    Gần triệt tiêu & 64 & 1.355 & 1.163 & 2.697 & 9,97 & 7,53 & 27,5 \\
    Gần triệt tiêu & 256 & 5.376 & 4.608 & 10.558 & 39,6 & 29,8 & 108 \\
    Gần triệt tiêu & 1.024 & 21.513 & 18.441 & 42.055 & 158 & 119 & 428 \\
    \bottomrule
  \end{tabular}
\end{table}

Trên mọi chuỗi $L$ phép MAC khác $0$ (với số vòng lặp bộ nhân $t_i = \mathrm{popcount}(X_{\mathrm{mant},i})$),
tổng số chu kỳ đo được tuân theo quan hệ đại số chính xác:
\[
  T_F = \sum_{i=1}^{L}\bigl(13 + \max(1,t_i)\bigr),\qquad
  T_{NF} = T_F + 3L,\qquad
  T_{\text{Quire}} = T_F + 23L + 62,
\]
trong đó \texttt{CLEAR} tốn $3$ chu kỳ, \texttt{READ} tốn $59$ chu kỳ và mỗi bước
tích lũy bank 36 bit của quire cộng thêm $23$ chu kỳ. Tuy nhiên, việc so sánh bằng
chu kỳ thuần túy chưa phản ánh đúng tốc độ thực tế vì ba lõi đạt tần số cực đại
sau định tuyến khác nhau. Khi quy đổi sang thời gian thực tại $F_{\max,\mathrm{R2R}}$
kiểm chứng sau định tuyến trên Artix-7 ($153{,}85$\,MHz ở $6{,}5$\,ns cho \texttt{MAC-F},
$135{,}14$\,MHz ở $7{,}4$\,ns cho \texttt{MAC-NF} và $98{,}16$\,MHz cho \texttt{MAC-Quire}
-- Bảng~\ref{tab:chain-cycles}), lõi \textbf{MAC-F} thực thi một chuỗi $1.024$ MAC mất
khoảng \textbf{120\,$\mu$s} ($119{,}9\text{--}120{,}1\,\mu$s), nhanh hơn \textbf{1,33 lần}
so với \texttt{MAC-NF} ($159{,}2\text{--}159{,}5\,\mu$s) và nhanh hơn \textbf{3,57 lần}
so với \texttt{MAC-Quire} ($428{,}4\text{--}428{,}8\,\mu$s).
Đổi lại, ở chuỗi gần triệt tiêu dài $1.024$ MAC, trung vị sai số tuyệt đối của
\texttt{MAC-Quire} giữ ở $1{,}73{\times}10^{-8}$ so với $2{,}17{\times}10^{-5}$
của hai lõi tích lũy theo mã Posit do quire loại bỏ hoàn toàn làm tròn trung gian
giữa các bước.

\section{Đánh giá trên bài toán ứng dụng thực tế}
\label{sec:app-benchmarks}
Để kiểm chứng giá trị thực tiễn ngoài tập dữ liệu ngẫu nhiên, các mô hình bit-exact
của \texttt{MAC-F}, \texttt{MAC-NF} và \texttt{MAC-Quire} được chạy trực tiếp trên
hai lớp bài toán xử lý tín hiệu và học máy gồm bộ lọc thông thấp FIR $32$ hệ số
($257$ mẫu đầu vào, $8.224$ phép MAC) và ba kịch bản tích vô hướng $128$ chiều
($200$ vectơ mỗi kịch bản), đặt trong đối chiếu trực tiếp với ba chuẩn dấu phẩy
động \texttt{IEEE-754 FP32}, \texttt{IEEE-754 FP16} và \texttt{BFloat16}
(\texttt{BF16}).

\begin{table}[!htbp]
  \centering
  \caption{Đánh giá trên hai bài toán ứng dụng thực tế: bộ lọc thông thấp FIR 32 hệ số ($257$ mẫu, $8.224$ phép MAC) và tích vô hướng $128$ chiều ($200$ vectơ mỗi kịch bản).}
  \label{tab:app-benchmarks}
  \footnotesize
  \setlength{\tabcolsep}{2.5pt}
  \renewcommand{\arraystretch}{1.12}
  \begin{tabularx}{\textwidth}{@{}>{\raggedright\arraybackslash}p{22.5mm} r r r r @{\hspace{6pt}}r @{\hspace{6pt}}r@{}}
    \toprule
    \textbf{Cấu hình lõi} & \textbf{FIR SNR} & \textbf{FIR RMSE} & \textbf{T.gian/mẫu} & \textbf{NN P95} & \textbf{DynScale P95} & \textbf{Cancel P95 (Lỗi)} \\
    & (dB) & & ($\mu$s tại $F_{\max}$) & ($|err|/|ref|$) & (Số lỗi Inf/UF) & (Số lỗi Inf/NaN) \\
    \midrule
    \textbf{IEEE-754 FP32} & $138{,}9$ & $6{,}00{\times}10^{-8}$ & --- & $1{,}32{\times}10^{-6}$ & \texttt{Inf} (\textbf{100} \texttt{Inf}) & \texttt{NaN} (\textbf{200/200}) \\
    \textbf{IEEE-754 FP16} & $62{,}2$ & $4{,}09{\times}10^{-4}$ & --- & $7{,}01{\times}10^{-3}$ & \texttt{Inf}/UF (\textbf{200/200}) & \texttt{NaN} (\textbf{200/200}) \\
    \textbf{BFloat16} & $44{,}1$ & $3{,}30{\times}10^{-3}$ & --- & $5{,}74{\times}10^{-2}$ & \texttt{Inf}/UF (\textbf{200/200}) & \texttt{NaN} (\textbf{200/200}) \\
    \textbf{MAC-F} ($n=12$) & $72{,}5$ & $1{,}25{\times}10^{-4}$ & $3{,}75$ & $5{,}66{\times}10^{-4}$ & $8{,}69{\times}10^{-3}$ (0) & mất phần dư ($0$ Inf) \\
    \textbf{MAC-F} ($n=8$) & $72{,}4$ & $1{,}28{\times}10^{-4}$ & $3{,}75$ & $5{,}58{\times}10^{-4}$ & $8{,}69{\times}10^{-3}$ (0) & mất phần dư ($0$ Inf) \\
    \textbf{MAC-NF} ($n=12$) & $72{,}5$ & $1{,}25{\times}10^{-4}$ & $4{,}98$ & $5{,}66{\times}10^{-4}$ & $8{,}27{\times}10^{-3}$ (0) & mất phần dư ($0$ Inf) \\
    \textbf{MAC-Quire} ($n=12$) & $72{,}5$ & $1{,}25{\times}10^{-4}$ & $14{,}0$ & $5{,}66{\times}10^{-4}$ & $1{,}14{\times}10^{-3}$ (0) & $\mathbf{7{,}73{\times}10^{-4}}$ (\textbf{0}) \\
    \bottomrule
  \end{tabularx}
  \vspace{2pt}
  {\footnotesize Ghi chú: Trong FIR và tích vô hướng chuẩn hóa \texttt{NN} ($128$ chiều), tỷ lệ cắt \texttt{approx\_cut} bằng $0\%$ ở $n=12$. Kịch bản \texttt{DynScale} phủ biên độ $10^{-45}\text{--}10^{43}$: \texttt{FP16} ($\max 65.504$) và \texttt{BF16} ($10^{\pm 38}$) hỏng $200/200$ vectơ ($100$ \texttt{Inf} + $100$ \texttt{underflow} về $0$), còn \texttt{FP32} hỏng $100/200$ \texttt{Inf} trên nửa lớn và chịu lượng tử hóa \textit{subnormal} nặng ở nửa nhỏ ($10^{-48}\text{--}10^{-42}$ tiệm cận biên $1{,}4{\times}10^{-45}$). \texttt{Cancel} cộng cặp $\pm 10^{45}$ kèm số hạng dư $10^{-44}$. Thời gian mỗi mẫu FIR tính từ chu kỳ đo ($577{,}0$, $673{,}0$ và $1.375$) chia $F_{\max,\mathrm{R2R}}$ kiểm chứng ($153{,}85$, $135{,}14$ và $98{,}16$\,MHz).}
\end{table}


Kết quả trong Bảng~\ref{tab:app-benchmarks} làm rõ phân vùng ưu thế và giới hạn
của từng kiến trúc:
\begin{itemize}
  \item \textbf{Xử lý tín hiệu số (FIR 32 hệ số) và tích vô hướng chuẩn hóa (\texttt{NN}):}
        Với cấu hình mặc định $n=12$, cả ba lõi Posit đạt tỷ lệ cắt \texttt{approx\_cut}
        bằng $0\%$, tái tạo đáp ứng bộ lọc FIR với $\mathrm{SNR} = 72{,}5$\,dB ($\mathrm{RMSE} = 1{,}25{\times}10^{-4}$)
        và sai số $\mathrm{P95}$ trên \texttt{NN} là $5{,}66{\times}10^{-4}$. Đặt cạnh
        các định dạng thực dụng trong suy luận biên, \texttt{MAC-F} vượt trội
        \textbf{$+10{,}3$\,dB} (sai số thấp hơn $3{,}27\times$) so với \texttt{FP16}
        ($62{,}2$\,dB) và \textbf{$+28{,}4$\,dB} so với \texttt{BF16} ($44{,}1$\,dB);
        trên \texttt{NN} 128 chiều, sai số $\mathrm{P95}$ của \texttt{MAC-F} thấp
        hơn \textbf{$12{,}4\times$} so với \texttt{FP16} ($7{,}01{\times}10^{-3}$)
        và \textbf{$101\times$} so với \texttt{BF16} ($5{,}74{\times}10^{-2}$).
        Đồng thời, cần ghi nhận khách quan rằng ba biến thể \texttt{MAC-F},
        \texttt{MAC-NF} và \texttt{MAC-Quire} cho cùng mức SNR/RMSE trên hai bài
        toán biên độ chuẩn này vì sai số cắt 12 bit của bộ nhân lặp chi phối tổng
        thể khi $|C|\sim |AB|$; ưu thế thực tế của \texttt{MAC-F} so với
        \texttt{MAC-NF} và \texttt{MAC-Quire} ở đây nằm ở thời gian thực thi mỗi
        mẫu ($3{,}75\,\mu$s so với $4{,}98$ và $14{,}0\,\mu$s), tiết kiệm 23,3\%
        LUT và giữ toàn vẹn hành vi làm tròn đơn bước ở biên bão hòa.
  \item \textbf{Tính toán dải động cực rộng (\texttt{DynScale}, $10^{-45}\text{--}10^{43}$):}
        Khi biên độ trải rộng ngoài ngưỡng chuẩn thông thường, \texttt{FP16} và
        \texttt{BFloat16} thất bại hoàn toàn $200/200$ vectơ ($100$ \texttt{Inf}
        và $100$ \texttt{underflow} về $0$), còn \texttt{FP32} tràn \texttt{Inf}
        trên $100/100$ vectơ nửa lớn và chịu lượng tử hóa \textit{subnormal} nặng
        ở nửa nhỏ sát ngưỡng $1{,}4{\times}10^{-45}$. Ngược lại, mọi lõi \texttt{Posit(32,3)}
        đều duy trì $200/200$ kết quả hữu hạn với sai số $\mathrm{P95}$ chỉ $0{,}869\%$
        ở \texttt{MAC-F} và $0{,}114\%$ ở \texttt{MAC-Quire}.
  \item \textbf{Triệt tiêu thảm khốc trên dải rộng (\texttt{Cancel}, $\pm 10^{45} + 10^{-44}$):}
        Cả \texttt{FP32}, \texttt{FP16} và \texttt{BF16} đều trả \texttt{NaN/Inf}
        trên $200/200$ vectơ; \texttt{MAC-F} và \texttt{MAC-NF} giữ được giá trị
        hữu hạn ($0$ \texttt{Inf}) nhưng mất phần dư $10^{-44}$ do cộng--làm tròn
        từng bước về mã 32 bit. Chỉ \texttt{MAC-Quire} bảo toàn phần dư trên lưới
        972 bit, đạt $\mathrm{P95} = 7{,}73{\times}10^{-4}$ với $0$ lỗi. Hai kịch
        bản \texttt{DynScale} và \texttt{Cancel} được dựng có chủ đích với biên
        độ nằm ngoài dải biểu diễn của \texttt{FP32} ngay từ đầu vào để minh họa
        dải động $2^{\pm 240}$; chúng không đại diện cho dữ liệu DSP hay mạng
        nơ-ron thông thường.
\end{itemize}

\section{Kết quả FPGA}
Các phép tổng hợp và định tuyến đều thực hiện ở chế độ OOC trên Vivado 2026.1
cho FPGA Artix-7 \texttt{xc7a35tcpg236-1} với chu kỳ clock mục tiêu 10\,ns (100\,MHz),
\texttt{uncertainty} 0,1\,ns, \texttt{input\_delay} 2,5\,ns và \texttt{output\_delay}
min/max 0/0,5\,ns; các wrapper đều giữ động cổng cấu hình \texttt{cfg\_ops} và
\texttt{cfg\_n}.

Trước hết, ở pha đối chứng tổng hợp kiến trúc cơ sở (Bảng~\ref{tab:nf-f-synthesis}
và Bảng~\ref{tab:f2f-stages} tại Phụ lục~\ref{app:fpga}), việc loại bỏ khối đóng gói và giải mã
tích trung gian giúp \textbf{MAC-F} giảm đáng kể tài nguyên và rút ngắn độ trễ
F2F lớn nhất từ 6,314\,ns xuống 5,430\,ns so với \textbf{MAC-NF}. Đối với
\textbf{MAC-Quire}, Bảng~\ref{tab:quire-synthesis} ghi nhận hiệu quả của việc
chuyển đổi kiến trúc backend từ dịch toàn vector 972 bit sang cập nhật cục bộ
theo từng từ 36 bit: tài nguyên LUT giảm từ 10.156 xuống 2.888 (giảm 71,6\%), FF
giảm từ 4.458 xuống 2.636 (giảm 40,9\%) và độ trễ F2F D/CE giảm gần một nửa từ
18,483\,ns xuống 9,579\,ns.

\begin{table}[htbp]
  \centering
  \caption{Đối chứng tổng hợp OOC kiến trúc cơ sở giữa MAC-NF và MAC-F trên Artix-7 A35T (100\,MHz).}
  \label{tab:nf-f-synthesis}
  \begin{tabular}{lrr}
    \hline
    Chỉ tiêu & MAC-NF & MAC-F \\
    \hline
    LUT & 1.582 & 1.235 \\
    FF & 916 & 781 \\
    F2F D/CE lớn nhất (ns) & 6,314 & 5,430 \\
    \hline
  \end{tabular}
\end{table}

\begin{table}[htbp]
  \centering
  \caption{Hiệu quả tối ưu backend Quire sau tổng hợp OOC trên Artix-7 A35T (100\,MHz).}
  \label{tab:quire-synthesis}
  \begin{tabular}{lrr}
    \hline
    Chỉ tiêu & Backend toàn vector & Backend theo từ 36 bit \\
    \hline
    LUT & 10.156 & 2.888 \\
    FF & 4.458 & 2.636 \\
    F2F D/CE lớn nhất (ns) & 18,483 & 9,579 \\
    \hline
  \end{tabular}
\end{table}

Để đánh giá trung thực hiệu năng sau định tuyến vật lý trên cùng một bộ ràng
buộc OOC 100\,MHz, Bảng~\ref{tab:fpga-synth-vs-route} đối chiếu trực tiếp hai pha
\textit{Synthesis} và \textit{Place \& Route} cho cả ba kiến trúc, đồng thời
trình bày bước nhảy định thời khi áp dụng kỹ thuật tách tải fanout tại tầng S1
trong cấu hình thiết kế hoàn thiện. Trong bảng, setup WNS toàn thiết kế được báo
cáo tách riêng khỏi setup WNS nội bộ \textit{reg-to-reg} trên miền \texttt{clk}.

\begin{table}[!htbp]
  \centering
  \caption{So sánh cùng ràng buộc OOC 100\,MHz (\texttt{xc7a35tcpg236-1}) giữa pha tổng hợp (\textit{Synth}) và định tuyến (\textit{Route}): tách riêng WNS toàn thiết kế (gồm I/O \texttt{in\_ready}) và WNS nội bộ \textit{reg-to-reg} trên miền \texttt{clk}}
  \label{tab:fpga-synth-vs-route}
  \footnotesize
  \setlength{\tabcolsep}{3.5pt}
  \begin{tabularx}{\textwidth}{@{}Y l r r r r r r@{}}
    \toprule
    \textbf{Cấu hình kiến trúc} & \textbf{Pha} & \textbf{LUT} & \textbf{FF} & \textbf{WNS toàn} & \textbf{WNS R2R} & \textbf{$F_{\max,\mathrm{R2R}}$} & \textbf{WHS (fail)} \\
     &  &  &  & \textbf{TK (ns)} & \textbf{(ns)} & \textbf{(MHz)} & \textbf{(ns)} \\
    \midrule
    MAC-NF (trước tách S1) & Synth & 1.582 & 916 & $+1{,}843$ & $+3{,}435$ & 152,32 & $-0{,}055$ (8) \\
    MAC-NF (trước tách S1) & Route & 1.558 & 916 & $+0{,}665$ & $+2{,}943$ & 141,70 & $+0{,}043$ (0) \\
    MAC-NF (tách S1) & Synth & 1.616 & 916 & $+1{,}843$ & $+3{,}435$ & 152,32 & $-0{,}055$ (8) \\
    MAC-NF (tách S1) & Route & 1.591 & 916 & $+0{,}669$ & $+2{,}639$ & 135,85 & $+0{,}041$ (0) \\
    \midrule
    MAC-F (trước tách S1) & Synth & 1.235 & 781 & $+2{,}689$ & $+4{,}319$ & 176,03 & $-0{,}055$ (7) \\
    MAC-F (trước tách S1) & Route & 1.211 & 781 & $+2{,}142$ & $+2{,}768$ & 138,27 & $+0{,}043$ (0) \\
    \textbf{MAC-F (tách S1, hoàn thiện)} & \textbf{Synth} & \textbf{1.244} & \textbf{781} & \textbf{$+2{,}689$} & \textbf{$+4{,}319$} & \textbf{176,03} & \textbf{$-0{,}055$ (7)} \\
    \textbf{MAC-F (tách S1, hoàn thiện)} & \textbf{Route} & \textbf{1.221} & \textbf{781} & \textbf{$+2{,}247$} & \textbf{$+3{,}524$} & \textbf{154,42} & \textbf{$+0{,}043$ (0)} \\
    \midrule
    MAC-Quire & Synth & 2.888 & 2.636 & $+0{,}170$ & $+0{,}170$ & 101,73 & $-0{,}055$ (7) \\
    MAC-Quire & Route & 2.873 & 2.636 & $-0{,}187$ & $-0{,}187$ & 98,16 & $+0{,}040$ (0) \\
    \bottomrule
  \end{tabularx}
  \vspace{2pt}
  {\footnotesize Ghi chú: Cả ba cấu hình có $\mathrm{DSP}=0$, $\mathrm{BRAM}=0$ ở cả hai pha; riêng MAC-Quire dùng thêm 174 F7~Mux và 1 F8~Mux. Các cấu hình NF/F đều đạt $\mathrm{TNS}=0{,}000$\,ns và 0 lỗi định tuyến ở pha Route; cấu hình \texttt{tách S1} giữ nguyên hoàn toàn datapath, số FF và độ trễ chu kỳ, chỉ tách 4 nhánh cho \texttt{choose\_a} trong \texttt{posit\_add\_prepare\_comb}. Đối với MAC-Quire, tổng hợp đạt $\mathrm{TNS}=0{,}000$\,ns nhưng sau định tuyến còn $\mathrm{TNS}=-19{,}336$\,ns với 280/5.271 điểm cuối intra-\texttt{clk} vi phạm (giá trị WHS trước route là ước lượng chưa đặt mạch, đều đóng về dương sau route). $F_{\max,\mathrm{R2R}} = 1000 / (10{,}000 - \mathrm{WNS}_{\mathrm{R2R}})$.}
\end{table}


Trên cấu hình trước khi tối ưu tầng S1, cả \texttt{MAC-NF} và \texttt{MAC-F} đều
đạt $\mathrm{TNS}=0$ và $\mathrm{THS}=0$ với $0$ lỗi định tuyến. Tuy nhiên, đường
setup găng nhất toàn thiết kế nằm trên nhánh tổ hợp từ thanh ghi đồng bộ đặt lại
ra cổng sẵn sàng nhận dữ liệu ($+0{,}665$\,ns ở \texttt{MAC-NF} và $+2{,}142$\,ns
ở \texttt{MAC-F}), vốn chịu ảnh hưởng trực tiếp từ giả định \texttt{output\_delay}
tại biên OOC. Khi truy vấn riêng miền nội bộ \textit{reg-to-reg}, \texttt{MAC-NF}
đạt $+2{,}943$\,ns ($F_{\max,\mathrm{R2R}} = 141{,}70$\,MHz) còn \texttt{MAC-F}
chỉ đạt $+2{,}768$\,ns ($F_{\max,\mathrm{R2R}} = 138{,}27$\,MHz) do tín hiệu điều
khiển hoán đổi toán hạng gánh tới $62$ nhánh tải ở tầng S1, làm khuếch đại trễ
định tuyến trên nhánh từ tầng chuẩn hóa sang bộ cộng.

Khi tách tín hiệu hoán đổi toán hạng thành bốn nhánh cân bằng tải tại khối chuẩn
bị S1 (giữ nguyên hoàn toàn datapath, số FF và độ trễ chu kỳ của \texttt{MAC-F}),
ở ràng buộc $10{,}000$\,ns, WNS nội bộ \textit{reg-to-reg} của \textbf{MAC-F} phục hồi
lên \textbf{$+3{,}524$\,ns} (ước lượng tuyến tính $F_{\max,\mathrm{R2R}} = 154{,}42$\,MHz,
tăng $+16{,}15$\,MHz) và WNS toàn thiết kế đạt \textbf{$+2{,}247$\,ns} (ước lượng
$F_{\max}$ toàn thiết kế $= 129{,}0$\,MHz) tại chi phí \textbf{1.221 LUT / 781 FF}
($\mathrm{WHS}=+0{,}043$\,ns). Trên cùng khối S1 dùng chung tại $10$\,ns, \texttt{MAC-NF}
đạt WNS toàn thiết kế $+0{,}669$\,ns (ước lượng $107{,}2$\,MHz; chưa đóng được ở
$8{,}000$\,ns với $\mathrm{WNS}=-0{,}074$\,ns) và nội bộ $+2{,}639$\,ns (ước lượng
$135{,}85$\,MHz) với 1.591 LUT / 916 FF. Để xác nhận thực chất thay vì dựa vào
ngoại suy tuyến tính đơn điểm, phép quét chu kỳ ràng buộc sau định tuyến với chiến
lược \texttt{Default} kiểm chứng rằng \textbf{MAC-F} đóng sạch định thời \textit{reg-to-reg}
đến \textbf{$T_{\mathrm{clk}}=6{,}500$\,ns} (\textbf{$F_{\max,\mathrm{R2R}} \ge 153{,}8$\,MHz},
cụ thể $153{,}85$\,MHz với $\mathrm{WNS}_{\mathrm{R2R}}=+0{,}632$\,ns,
$\mathrm{WHS}_{\mathrm{R2R}}=+0{,}043$\,ns, $1.225$ LUT / $781$ FF) và đóng sạch toàn
thiết kế đến \textbf{$T_{\mathrm{clk}}=7{,}000$\,ns} (\textbf{$F_{\max}$ toàn thiết kế
$\ge 142{,}9$\,MHz}, cụ thể $142{,}86$\,MHz với $\mathrm{WNS}=+0{,}045$\,ns dưới
ngân sách I/O $2{,}5/0{,}5$\,ns), trong khi \texttt{MAC-NF} đóng \textit{reg-to-reg} đến
\textbf{$T_{\mathrm{clk}}=7{,}400$\,ns} (\textbf{$F_{\max,\mathrm{R2R}} \ge 135{,}1$\,MHz},
cụ thể $135{,}14$\,MHz với $\mathrm{WNS}_{\mathrm{R2R}}=+0{,}209$\,ns,
$\mathrm{WHS}_{\mathrm{R2R}}=+0{,}046$\,ns, $1.599$ LUT / $916$ FF). Ở mức siết
ràng buộc $T_{\mathrm{clk}}=6{,}500$\,ns, giá trị $\mathrm{WNS}=-0{,}513$\,ns
trên tổng các đường chỉ xuất phát từ việc giữ ràng buộc \texttt{set\_output\_delay}
ngoại vi cố định mức $100$\,MHz trên cổng \texttt{in\_ready}; khi bỏ trễ I/O ngoại vi
(chế độ đo lõi nội bộ), \texttt{MAC-F} đóng sạch toàn bộ đường ở $6{,}500$\,ns với
$\mathrm{WNS}=\mathrm{WNS}_{\mathrm{R2R}}=+0{,}182$\,ns ($\mathrm{WHS}=+0{,}035$\,ns,
$1.229$ LUT). Đồng thời, khi quét ba chiến lược đặt--định tuyến (\texttt{Default},
\texttt{Explore}, \texttt{ExtraTimingOpt}) tại $10$\,ns, $\mathrm{WNS}_{\mathrm{R2R}}$
dao động trong $[+3{,}184; +3{,}524]$\,ns đối với \texttt{MAC-F} và $[+2{,}639; +2{,}994]$\,ns
đối với \texttt{MAC-NF}, khẳng định chênh lệch $+18{,}71$\,MHz của \texttt{MAC-F} là
ổn định.
Ngược lại, \texttt{MAC-Quire} đạt $+0{,}170$\,ns ở pha tổng hợp nhưng chuyển sang
$-0{,}187$\,ns sau định tuyến ($\mathrm{TNS}=-19{,}336$\,ns, 280/5.271 điểm cuối
intra-\texttt{clk}, $F_{\max,\mathrm{R2R}}=98{,}16$\,MHz) do đường găng nằm sâu
trong mạch giải mã chọn từ của bank tích lũy (chi tiết từng đường găng thanh ghi
được liệt kê trong Bảng~\ref{tab:route-critical}).

\subsection{Thông lượng, công suất và đối chiếu với Xilinx FP32 cùng công trình liên quan}
Để hoàn thiện bức tranh đánh giá hiệu năng phần cứng, đồ án trích xuất trực tiếp
\texttt{Vivado report\_power} (100\,MHz vectorless trên \texttt{routed.dcp})
và triển khai thêm ba cấu hình chuẩn đối chứng dùng IP \texttt{Xilinx Floating-Point v7.1}
\texttt{FP32}~\cite{amd_pg060_fp} trên cùng chip \texttt{xc7a35tcpg236-1} với cùng
ràng buộc OOC 100\,MHz: \texttt{FMA Full\_Usage} ($4$ DSP, trễ $20$ chu kỳ),
\texttt{FMA Medium\_Usage} ($2$ DSP, trễ $17$ chu kỳ) và biến thể $0$-DSP phải
ghép nối tiếp khối nhân logic với khối cộng trừ (\texttt{Multiply No\_Usage + Add\_Subtract},
trễ $22$ chu kỳ, vì lõi FMA của Xilinx không cho phép cấu hình \texttt{No\_Usage}).

\begin{table}[!htbp]
  \centering
  \caption{Tổng hợp hiệu năng, thông lượng, công suất (\texttt{Vivado report\_power} tại 100\,MHz) và chỉ số hiệu quả trên \texttt{xc7a35tcpg236-1} (3 chữ số có nghĩa).}
  \label{tab:fpga-power-throughput}
  \footnotesize
  \setlength{\tabcolsep}{2.8pt}
  \renewcommand{\arraystretch}{1.12}
  \begin{tabularx}{\textwidth}{@{}>{\raggedright\arraybackslash}p{31mm} r r r r r r r r@{}}
    \toprule
    \textbf{Kiến trúc lõi} & \textbf{LUT} & \textbf{FF} & \textbf{DSP} & \textbf{$F_{\max,\mathrm{R2R}}$} & \textbf{Chu kỳ} & \textbf{Thông lượng} & \textbf{$P_{\mathrm{dyn}}/P_{\mathrm{tot}}$} & \textbf{$E_{\mathrm{dyn}}/\mathrm{MAC}$} \\
    & & & & (MHz) & (II/Lat) & (M\,MAC/s) & (mW) & (nJ) \\
    \midrule
    \textbf{MAC-F} (tách S1) & \textbf{1.221} & \textbf{781} & \textbf{0} & 153,8 & 14--25 (18,0) & 6,15--11,0 (8,53) & \textbf{4 / 72} & 0,721 \\
    \textbf{MAC-NF} (tách S1) & 1.591 & 916 & 0 & 135,1 & 17--28 (21,0) & 4,83--7,95 (6,43) & 12 / 80 & 2,52 \\
    \textbf{MAC-Quire} (mỗi MAC) & 2.873 & 2.636 & 0 & 98,2 & 37--48 (41,0) & 2,05--2,65 (2,39) & 29 / 98 & 11,9 \\
    \midrule
    \textbf{Xilinx FP32 FMA Full} & 765 & 1.318 & 4 & 180 & 1 / 20 & 180 (9,00*) & 25 / 94 & 0,250 \\
    \textbf{Xilinx FP32 FMA Med} & 777 & 1.285 & 2 & 191 & 1 / 17 & 191 (11,2*) & 21 / 89 & 0,210 \\
    \textbf{Xilinx FP32 0-DSP} & 993 & 1.520 & 0 & 195 & 1 / 22 & 195 (8,88*) & 22 / 91 & 0,220 \\
    \bottomrule
  \end{tabularx}
  \vspace{2pt}
  {\footnotesize Ghi chú: Ba lõi Posit có $II \in [14,25]$, $[17,28]$ và $[37,48]$ theo $0 \le t \le 12$ vòng nhân (trung bình trong ngoặc ứng với $5{,}03$ vòng ở \texttt{OPS1/n12}). Tần số $F_{\max,\mathrm{R2R}}$ và thông lượng của MAC-F/MAC-NF đã được route thực tế tại chu kỳ mục tiêu ($153{,}85$\,MHz tại $6{,}5$\,ns và $135{,}14$\,MHz tại $7{,}4$\,ns), còn $F_{\max}$ của các lõi \texttt{Xilinx floating\_point:7.1} ($180$/$191$/$195$\,MHz) được suy ra từ $\mathrm{WNS}_{\mathrm{R2R}}$ ở lượt $10$\,ns ($1000/(10-\mathrm{WNS}_{\mathrm{R2R}})$), chưa route thực ở chu kỳ tương ứng. Tài nguyên LUT/FF và công suất \texttt{report\_power} không vectơ ($\pm 0{,}5$\,mW) giữ nguyên ở điểm cơ sở $100$\,MHz. Với các lõi Xilinx đường ống $II=1$, thông lượng luồng độc lập đạt $F_{\max}$ và $E_{\mathrm{dyn}} = P_{\mathrm{dyn}}\cdot 10\,\text{ns} = 0{,}210\text{--}0{,}250$\,nJ (chưa tính chu kỳ chờ nếu tích lũy hồi tiếp không \textit{clock-gate}), còn thông lượng có dấu * ($F_{\max}/\mathrm{Lat}$) áp dụng cho chuỗi tích lũy phụ thuộc đơn lẻ $C_{i+1}=D_i$. Chỉ số diện tích--thời gian $\mathrm{LUT}\times\mathrm{ns}/\mathrm{MAC}$ tại $F_{\max}$ lần lượt là $1{,}43{\times}10^5$ (\texttt{MAC-F}), $2{,}48{\times}10^5$ (\texttt{MAC-NF}) và $1{,}20{\times}10^6$ (\texttt{MAC-Quire}), so với $4{,}07{\times}10^3\text{--}5{,}09{\times}10^3$ ở luồng độc lập $II=1$ của Xilinx ($8{,}50{\times}10^4\text{--}1{,}12{\times}10^5$ khi chờ hồi tiếp).}
\end{table}


Bảng~\ref{tab:fpga-power-throughput} làm rõ sự đánh đổi kiến trúc giữa lõi tuần
tự tiết kiệm tài nguyên và lõi đường ống chuẩn IEEE-754:
\begin{itemize}
  \item \textbf{Thông lượng và hiệu quả năng lượng nội bộ:} Nhờ rút 3 chu kỳ/MAC
        và nâng $F_{\max,\mathrm{R2R}}$ kiểm chứng lên $\ge 153{,}8$\,MHz ($153{,}85$\,MHz),
        \textbf{MAC-F} đạt thông lượng $6{,}15\text{--}11{,}0$\,M\,MAC/s (trung bình
        $8{,}53$\,M\,MAC/s ở $5{,}03$ vòng lặp), cao hơn 32,8\% so với \texttt{MAC-NF}
        ($4{,}83\text{--}7{,}95$\,M\,MAC/s, TB $6{,}43$) và gấp $3{,}57\times$ so với
        \texttt{MAC-Quire} ($2{,}05\text{--}2{,}65$\,M\,MAC/s, TB $2{,}39$). Đồng thời, ở phép
        ước lượng \textit{vectorless} tại $100$\,MHz
        (độ phân giải báo cáo $\pm 0{,}5$\,mW), công suất động của \texttt{MAC-F} là
        \textbf{4\,mW} (tổng $72$\,mW gồm $68$\,mW tĩnh nền chip), thấp hơn
        \texttt{MAC-NF} ($12$\,mW) và giảm năng lượng động mỗi phép MAC về
        \textbf{$0{,}721$\,nJ} (so với $2{,}52$\,nJ của \texttt{MAC-NF} và $11{,}9$\,nJ
        của \texttt{MAC-Quire}).
  \item \textbf{So sánh khách quan với IP \texttt{Xilinx FP32}:} Khi vận hành dạng
        luồng độc lập ($II=1$), các lõi đường ống \texttt{Xilinx FP32} (với
        $F_{\max}$ $180\text{--}195$\,MHz suy ra tuyến tính từ lượt $10$\,ns, chưa
        route thực ở chu kỳ tương ứng) vượt trội về thông lượng
        ($180\text{--}195$\,M\,MAC/s), dùng ít LUT hơn ($765\text{--}993$ so với
        $1.221$ LUT của \texttt{MAC-F}), đạt chỉ số $\mathrm{LUT}\times\mathrm{ns}/\mathrm{MAC}$
        tốt hơn ($4{,}07{\times}10^3\text{--}5{,}09{\times}10^3$ so với
        $1{,}43{\times}10^5$), tiêu tốn năng lượng động/MAC thấp hơn khoảng $3\times$
        ($0{,}210\text{--}0{,}250$\,nJ so với $0{,}721$\,nJ do thời gian tính 1 chu kỳ
        thay vì $18{,}0$ chu kỳ) và đạt độ chính xác FIR cao hơn rõ rệt
        ($138{,}9$\,dB so với $72{,}5$\,dB); đây là cái giá đánh đổi cố hữu của
        thiết kế lặp tuần tự và xấp xỉ 12 bit. Đổi lại, \texttt{Xilinx FP32} tiêu
        tốn tới \textbf{1.285--1.520 FF} (gấp $1{,}65\text{--}1{,}95\times$ mức $781$ FF
        của \texttt{MAC-F}), cần $2\text{--}4$ khối \texttt{DSP48E1} ở biến thể FMA
        và tiêu thụ công suất động \textbf{21--25\,mW} (khoảng $5\text{--}7\times$
        của \texttt{MAC-F}, tính cả độ phân giải $\pm 0{,}5$\,mW của \texttt{report\_power}).
        Trong kịch bản chỉ có \textbf{một chuỗi tích lũy phụ thuộc đơn lẻ}
        $C_{i+1}=D_i$ (nơi mỗi phép mới phải chờ hết $17\text{--}22$ chu kỳ trễ
        đường ống), thông lượng thực dụng của \texttt{Xilinx FP32 0-DSP} giảm về
        $195{,}3 / 22 = 8{,}88$\,M\,MAC/s, ngang bằng $8{,}53$\,M\,MAC/s của
        \texttt{MAC-F} trong khi \texttt{MAC-F} dùng ít hơn $48{,}6\%$ FF, giữ
        \texttt{0 DSP} và mở rộng dải động lên $2^{\pm 240}$; ngược lại, nếu xen
        kẽ nhiều chuỗi độc lập (ví dụ nhiều đầu ra FIR song song), lõi
        \texttt{Xilinx FP32} lấy lại nhịp $II=1$ và vượt trội về thông lượng.
\end{itemize}

\begin{table}[!htbp]
  \centering
  \caption{Đối chiếu định lượng và kiến trúc với IP Xilinx \texttt{FP32} (trên cùng \texttt{xc7a35t}) và các nghiên cứu Posit 32 bit đã công bố.}
  \label{tab:literature-comparison}
  \footnotesize
  \setlength{\tabcolsep}{2.8pt}
  \renewcommand{\arraystretch}{1.04}
  \begin{tabularx}{\textwidth}{@{}>{\raggedright\arraybackslash}p{29mm} >{\raggedright\arraybackslash}p{24mm} r r r r Y@{}}
    \toprule
    \textbf{Thiết kế / Công trình} & \textbf{Phép toán / Chuẩn} & \textbf{Nền tảng} & \textbf{LUT} & \textbf{FF} & \textbf{DSP} & \textbf{Tần số \& Đặc điểm chính} \\
    \midrule
    \textbf{Đồ án: MAC-F} & Fused MAC \texttt{Posit(32,3)} & Artix-7 A35T & 1.221 & 781 & 0 & $\ge 153{,}8$\,MHz (R2R) / $\ge 142{,}9$\,MHz (TK, route 7,0\,ns; ước lượng 10\,ns: 129,0\,MHz); 14--25 cyc; $P_{\mathrm{dyn}}=4$\,mW. \\
    \textbf{Đồ án: MAC-NF} & Non-fused \texttt{Posit(32,3)} & Artix-7 A35T & 1.591 & 916 & 0 & $\ge 135{,}1$\,MHz (R2R) / 100\,MHz (TK, ước lượng 10\,ns: 107,2\,MHz; lỗi 8,0\,ns); 17--28 cyc; $P_{\mathrm{dyn}}=12$\,mW. \\
    \textbf{Đồ án: MAC-Quire} & Quire MAC \texttt{Posit(32,3)} & Artix-7 A35T & 2.873 & 2.636 & 0 & 98,2\,MHz; lưới 972 bit theo từ 36 bit; $P_{\mathrm{dyn}}=29$\,mW. \\
    \midrule
    Xilinx FP32 FMA Med~\cite{amd_pg060_fp} & Fused FMA \texttt{FP32} & Artix-7 A35T & 777 & 1.285 & 2 & 191\,MHz (ước lượng từ 10\,ns); đường ống $II=1$, Lat=17; $P_{\mathrm{dyn}}=21$\,mW. \\
    Xilinx FP32 0-DSP~\cite{amd_pg060_fp} & Mul+Add \texttt{FP32} & Artix-7 A35T & 993 & 1.520 & 0 & 195\,MHz (ước lượng từ 10\,ns); đường ống $II=1$, Lat=22; $P_{\mathrm{dyn}}=22$\,mW. \\
    \midrule
    Jaiswal \& So (PACoGen)~\cite{jaiswal2019pacogen} & Mul / Add \texttt{Posit(32,es)} & Virtex-7 & 854 / 946 & --- & 1 / 0 & 227 / 244\,MHz; đường ống tách rời Mul ($4$ cyc) và Add ($5$ cyc), $II=1$. \\
    Murillo et al. (FloPoCo)~\cite{murillo2020flopoco,murillo2022energy} & Gen Posit \texttt{Posit(32,es)} & Virtex-7 / 7-s & 572--682 & --- & 4 & Nhân chính xác dùng 4 DSP (mốc tham chiếu tổng hợp bởi~\cite{murillo2021plam}). \\
    Murillo et al. (PLAM)~\cite{murillo2021plam} & Approx Mul \texttt{Posit(32,es)} & FPGA 7-series & 528 & --- & 0 & Nhân xấp xỉ logarit $0$ DSP (đơn vị nhân đơn lẻ, chưa gồm Add/Quire). \\
    Mallasén et al. (PERCIVAL)~\cite{mallasen2022percival} & PAU+Quire \texttt{Posit(32,2)} & Kintex-7 325T & 11.879 & 2.985 & 4 & Khối số học tích hợp FPU/Quire đầy đủ trong nhân RISC-V (khác phạm vi IP MAC). \\
    \midrule
    Chaurasiya et al.~\cite{chaurasiya2018parameterized} & Gen Mul/Add \texttt{Posit(32,x)} & FPGA / ASIC & --- & --- & --- & Bộ sinh tham số hóa sớm; ghép riêng Add và Mul không hợp nhất lặp. \\
    Zhang \& Ko~\cite{zhang2020positmul} & Mul \texttt{Posit(32,x)} & FPGA / ASIC & --- & --- & --- & Phân vùng sub-multiplier $15{\times}15$ ngắt động, giảm $\approx 16\%$ công suất nhân. \\
    Norris \& Kim~\cite{norris2021iterativeposit} & Iterative Mul \texttt{Posit(32,3)} & Xilinx FPGA & --- & --- & --- & Nền tảng nhân lặp xấp xỉ 0 DSP (OPS/SAC/SBM); chưa tích hợp bộ cộng MAC hay Quire. \\
    \bottomrule
  \end{tabularx}
  \vspace{2pt}
  {\footnotesize Ghi chú: Do khác biệt về dòng chip (Artix-7 tốc độ $-1$ so với Kintex-7/Virtex-7), độ rộng phép toán (IP MAC tuần tự so với đường ống $II=1$ hoặc RISC-V PAU) và công cụ tổng hợp, bảng chỉ dùng để đối chiếu định vị kiến trúc; riêng 5 dòng đầu đo trực tiếp trên cùng quy trình OOC \texttt{xc7a35tcpg236-1} của Vivado 2026.1. Dấu ``---'' biểu thị tài liệu không tách riêng FF cho từng đơn vị con.}
\end{table}


Đặt trong bối cảnh nghiên cứu quốc tế (Bảng~\ref{tab:literature-comparison}), do
khác biệt dòng FPGA, công nghệ tổng hợp và độ rộng phép toán ($II=1$ đơn lẻ hay
hệ MAC tuần tự), số liệu công trình bên ngoài được dùng để định vị kiến trúc thay
vì so trực tiếp như cùng ràng buộc: khác với các bộ sinh phép toán tách rời như
\texttt{PACoGen}~\cite{jaiswal2019pacogen} ($854$ LUT + $1$ DSP cho Mul và $946$ LUT
cho Add trên Virtex-7), bộ nhân logarit $0$-DSP đơn lẻ \texttt{PLAM}~\cite{murillo2021plam}
($528$ LUT, tham chiếu bởi \texttt{FloPoCo} $4$-DSP~\cite{murillo2020flopoco,murillo2022energy}),
hay khối PAU+Quire đầy đủ trong nhân RISC-V \texttt{PERCIVAL}~\cite{mallasen2022percival}
($11.879$ LUT / $4$ DSP trên Kintex-7), lõi \texttt{MAC-F} đóng gói trọn vẹn chuỗi
Giải mã--Nhân lặp--Cộng 5 tầng--Packer RNE với \textbf{0 DSP} và \textbf{1.221 LUT}.

\section{Phân tích độ nhạy của phép đo sai số}

Phân vị 95 mô tả giá trị tại một vị trí trong phân bố,
còn max mô tả mẫu bất lợi nhất trong corpus. Chúng
phản ứng khác nhau khi một nhóm nhỏ giao dịch giữ
thêm số hạng. Từ n8 đến n9, số cặp lấy mẫu đều còn
term giảm từ 554 xuống 43; tới n10 bằng 0 trong nhóm
này. P95 gần với mức sai số do cắt đầu vào ngay từ
n8, nhưng max vẫn thay đổi theo những mẫu hiếm còn term.
Trong nhóm chủ đích, có cặp cần đủ 12 vòng. Vì vậy
không thể suy rằng mọi dữ liệu đều đủ ở n10.

\begin{engineeringtable}{Các mốc ở đuôi phân bố trên nhóm lấy đều, OPS1 (3 chữ số có nghĩa).}{tab:tail-sensitivity}
\begin{tabular}{rrrr}
\toprule
n & Cặp còn term & P95 sai số tổng (\%) & Max quan sát (\%) \\
\midrule
7 & 3.724 & 0,0320 & 0,351 \\
8 & 554 & 0,0299 & 0,156 \\
9 & 43 & 0,0297 & 0,0582 \\
10 & 0 & 0,0297 & 0,0460 \\
12 & 0 & 0,0297 & 0,0460 \\
\bottomrule
\end{tabular}
\end{engineeringtable}

Phân tầng regime giúp quan sát rõ các miền ít xuất hiện khi lấy đều theo mã. Với
$n=12$, nhóm phân tầng đạt $\mathrm{P95} = 0{,}0218\%$, giá trị cực đại $0{,}0438\%$
và trung bình $2{,}66$ vòng lặp (so với trung bình $5{,}03$ vòng ở nhóm lấy đều).
Sự khác biệt này phản ánh mật độ bit $1$ và độ dài fraction khác nhau giữa các dải
\textit{regime}; do đó số liệu của nhóm phân tầng được theo dõi tách biệt khỏi phân
bố ứng dụng thực tế.

\section{Từ sai số tích đến sai số MAC}

Sai số cắt fraction được bảo toàn giống nhau qua \texttt{MAC-NF} và \texttt{MAC-F}
cho cùng một tích đầu vào. Khi toán hạng cộng $C$ cùng dấu với $P$, độ lớn tổng tăng
nên cùng một sai lệch tích thường chiếm tỷ lệ tương đối nhỏ hơn; ngược lại, khi $C$
trái dấu, mẫu số thu hẹp và bước lượng tử hóa trung gian $R(\widehat P)$ của
\texttt{MAC-NF} làm xuất hiện sai khác phân kỳ so với tổng đầy đủ.

Trên $108.703$ ca gần $0$ ($r \ne 0$) thuộc nhóm ép triệt tiêu $C = R(-P)$, hai
chỉ số bất biến theo thang trong Bảng~\ref{tab:cancel-target-metrics} phản ánh bản
chất sai số chính xác hơn nhiều so với sai số tuyệt đối thô vốn biến thiên hàng
chục bậc theo hệ số mũ của \texttt{Posit(32,3)}. Đồng thời, vì $C$ bị lượng tử hóa
trước khi đưa vào RTL, một số cặp chạm ngưỡng kẹp hoặc bước nhảy thưa khiến $P+C$
không còn thỏa điều kiện gần $0$; việc phân lớp lại theo tổng chính xác $r = P + C$
trước khi thống kê giúp tách bạch hoàn toàn nhóm $r = 0$ khỏi các ca kẹp ngưỡng
ở lớp \texttt{regular}.

\section{Ý nghĩa của kết quả chuỗi}

Trong vận hành theo chuỗi, \texttt{MAC-NF} và \texttt{MAC-F} cập nhật trạng thái
\texttt{Posit(32,3)} sau mỗi phép hoàn tất nên sai khác làm tròn trung gian tích tụ
qua từng bước hồi tiếp. Trái lại, \texttt{MAC-Quire} tích lũy trực tiếp trên lưới
972 bit và chỉ làm tròn duy nhất tại lệnh \texttt{READ}. Ở nhóm chuỗi gần triệt tiêu
có nhiễu nhỏ tại bit thấp bị cắt của cặp toán hạng cuối, các tích đã cắt gọn triệt
tiêu lẫn nhau giúp \texttt{MAC-Quire} giữ sai số ở bậc $10^{-8}$, tránh lỗi tích tụ
trạng thái giữa chuỗi.

Tại độ dài $1.024$ MAC, độ trễ tuyệt đối theo xung nhịp thực (Bảng~\ref{tab:chain-cycles})
giúp phân định rõ ranh giới lựa chọn: khi đặt trong hệ thống xử lý thời gian thực
cần tần số cao và tài nguyên gọn, \texttt{MAC-F} nhanh hơn đáng kể so với cả hai
biến thể còn lại dù hai lõi \texttt{MAC-NF}/\texttt{MAC-F} có sai số chuỗi ngang nhau.

\begin{engineeringtable}{Sai số tuyệt đối cuối của nhóm gần triệt tiêu ($4$ chuỗi/lõi, 3 chữ số có nghĩa).}{tab:near-chain-detail}
\begin{tabular}{rrrrr}
\toprule
MAC/chuỗi & Trung vị NF/F & Max NF/F & Trung vị quire & Max quire \\
\midrule
16 & $6{,}94{\times}10^{-8}$ & $9{,}97{\times}10^{-8}$ & $9{,}47{\times}10^{-9}$ & $1{,}34{\times}10^{-8}$ \\
64 & $1{,}54{\times}10^{-7}$ & $7{,}07{\times}10^{-7}$ & $9{,}82{\times}10^{-9}$ & $1{,}02{\times}10^{-8}$ \\
256 & $2{,}14{\times}10^{-6}$ & $5{,}59{\times}10^{-6}$ & $8{,}30{\times}10^{-9}$ & $9{,}36{\times}10^{-9}$ \\
1.024 & $2{,}17{\times}10^{-5}$ & $1{,}28{\times}10^{-4}$ & $1{,}73{\times}10^{-8}$ & $2{,}61{\times}10^{-8}$ \\
\bottomrule
\end{tabular}
\end{engineeringtable}

Bảng~\ref{tab:near-chain-detail} dùng giá trị cuối so với
tổng tích đủ fraction, trên nhóm có các cặp tích đối dấu
và một perturbation ở bit thấp của cặp cuối. Với bốn mẫu
mỗi độ dài, các hàng mô tả giá trị đã quan sát. Chúng không
đủ để suy một luật tăng sai số theo chiều dài cho mọi dãy.
Giá trị trung vị của NF và F bằng nhau trong nhóm này,
dù số chu kỳ khác nhau. Vì vậy lựa chọn theo tài nguyên,
chu kỳ và độ chính xác phải giữ ba đại lượng tách biệt.

\section{Phân tích đường găng và tối ưu hóa vật lý sau định tuyến}

\subsection{Phân biệt tài nguyên giữa tổng hợp và định tuyến}

Ở pha tổng hợp cơ sở (Bảng~\ref{tab:nf-f-synthesis}), \texttt{MAC-F} dùng ít
hơn 21,9\% LUT (347 LUT) và 14,7\% FF (135 FF) so với \texttt{MAC-NF}. Chênh
lệch FF xuất phát trực tiếp từ ngân hàng thanh ghi của khối đóng gói tích trung
gian mà \texttt{MAC-F} loại bỏ, còn tỷ lệ LUT phụ thuộc thêm vào chiến lược gộp
logic điều khiển và ánh xạ chuỗi cộng của Vivado qua từng pha triển khai.

Sau định tuyến OOC 100\,MHz, cấu hình trước tách tải fanout của \texttt{MAC-F}
dùng 1.211 LUT / 781 FF (so với 1.558 LUT / 916 FF của \texttt{MAC-NF}). Khi
tách bốn nhánh cho tín hiệu hoán đổi toán hạng tại tầng chuẩn bị S1, thiết kế
hoàn thiện của \textbf{MAC-F} sử dụng 1.221 LUT (chỉ tăng $+10$ LUT) và giữ nguyên
781 FF, trong khi \texttt{MAC-NF} dùng 1.591 LUT và 916 FF; như vậy \textbf{MAC-F}
tiếp tục tiết kiệm 370 LUT (\textbf{23,3\%}) và 135 FF (\textbf{14,7\%}) so với \texttt{MAC-NF}.

\subsection{Phân tích sụt giảm timing hậu route và hiệu quả tách fanout tầng S1}

\begin{engineeringtable}{Các đường găng quyết định slack sau định tuyến OOC 100\,MHz: tách riêng đường cổng \texttt{in\_ready} và đường \textit{reg-to-reg} nội bộ trước và sau tối ưu tầng S1.}{tab:route-critical}
\begin{tabularx}{\linewidth}{@{}>{\raggedright\arraybackslash}p{34mm}>{\raggedright\arraybackslash}p{15mm}Y Y@{}}
\toprule
Lõi (trạng thái) & Slack (ns) & Điểm \textit{launch} & Điểm \textit{capture} \\
\midrule
NF setup (I/O, trước tách) & $+0{,}665$ & Packer cuối, \texttt{meta\_sync2\_reg} & Cổng \texttt{in\_ready}. \\
NF setup (R2R, trước tách) & $+2{,}943$ & \texttt{u\_product\_pack/\allowbreak d\_q\_reg[15]} & \texttt{product\_frac\_\allowbreak q\_reg[5]} (Parser tích). \\
NF setup (I/O, tách S1) & $+0{,}669$ & \texttt{u\_pack/p1\_valid\_\allowbreak q\_reg} & Cổng \texttt{in\_ready}. \\
NF setup (R2R, tách S1) & $+2{,}639$ & \texttt{u\_core/acc\_q\_reg[20]} & \texttt{u\_product\_pack/\allowbreak p1\_guard\_bit\_q\_reg} (\texttt{dp=7,256\,ns}). \\
F setup (I/O, trước tách) & $+2{,}142$ & Wrapper, \texttt{meta\_sync2\_reg} & Cổng \texttt{in\_ready}. \\
F setup (R2R, trước tách) & $+2{,}768$ & \texttt{core\_sf\_\allowbreak inc\_q\_reg[4]} & \texttt{u\_add/\allowbreak s1\_sf\_q\_reg[1]} (Adder S1). \\
\textbf{F setup (I/O, hoàn thiện)} & \textbf{$+2{,}247$} & \texttt{u\_b/p2\_valid\_q\_reg} & Cổng \texttt{in\_ready}. \\
\textbf{F setup (R2R, hoàn thiện)} & \textbf{$+3{,}524$} & \texttt{c\_frac\_q\_reg[9]} & \texttt{u\_add/\allowbreak s1\_x\_frac\_q\_reg[1]} (\texttt{dp=6,321\,ns}). \\
Quire setup (R2R) & $-0{,}187$ & \texttt{u\_bank/\allowbreak offset\_q\_reg[4]} & \texttt{g\_words[16].\allowbreak q\_words\_reg[16][8]/CE}. \\
NF hold (R2R, tách S1) & $+0{,}041$ & \texttt{u\_product\_pack/\allowbreak p1\_special\_sel\_q\_reg[2]} & \texttt{u\_product\_pack/\allowbreak d\_q\_reg[27]}. \\
F hold (R2R, hoàn thiện) & $+0{,}043$ & Adder, \texttt{flags\_out\_reg[0]} & Packer P1, \texttt{p1\_flags\_q\_reg[0]}. \\
Quire hold (R2R) & $+0{,}040$ & \texttt{pending\_words\_\allowbreak reg[11][32]} & \texttt{q\_words\_\allowbreak reg[11][32]}. \\
\bottomrule
\end{tabularx}
\end{engineeringtable}

Độ trễ F2F trong Bảng~\ref{tab:f2f-stages} (Phụ lục~\ref{app:fpga}) phản ánh
thời gian lan truyền tổ hợp lớn nhất tới từng nhóm thanh ghi đích sau tổng hợp,
dùng để nhận diện sớm tầng cần theo dõi. WNS toàn thiết kế sau định tuyến bao
gồm thêm độ lệch clock và quỹ định thời đường cổng ngoài (chi tiết từng thanh ghi
được ghi nhận tại Bảng~\ref{tab:route-critical}). Dữ liệu hậu \textit{route}
làm rõ hiện tượng đảo chiều định thời ở cấu hình trước tối ưu và hiệu quả của kỹ
thuật tách tải:
\begin{itemize}
  \item Ở pha tổng hợp, \texttt{MAC-F} vượt trội hoàn toàn với WNS \textit{reg-to-reg}
        đạt $+4{,}319$\,ns (5 tầng logic) so với $+3{,}435$\,ns của \texttt{MAC-NF}
        (6 tầng logic).
  \item Tuy nhiên sau định tuyến, đường găng nội bộ \textit{reg-to-reg} của
        \texttt{MAC-F} ban đầu sụt từ $+4{,}319$\,ns xuống còn $+2{,}768$\,ns
        ($138{,}27$\,MHz, trễ đường dữ liệu $7{,}018$\,ns với \textbf{73,67\% là
        trễ định tuyến}). Nguyên nhân vật lý nằm ở nhánh chuẩn hóa từ thanh ghi
        số mũ tích đi vào khối tiền xử lý cộng, nơi tín hiệu chọn toán hạng có
        độ phân tải tới $62$ nhánh song song trên toàn độ rộng 26 bit fraction và
        metadata tầng S1 (cùng nút tải $130$ nhánh ở bộ tích lũy), gây nghẽn
        định tuyến cục bộ.
\end{itemize}

Giải pháp nhân bản bốn nhánh cân bằng tải tại khối chuẩn bị S1 (giữ nguyên toàn
bộ datapath, số FF và độ trễ chu kỳ) đã dịch chuyển đường găng nội bộ sang nhánh
căn chỉnh thanh ghi cộng với tỷ lệ trễ cân đối hơn (38,78\% logic / 61,22\% định
tuyến), nâng WNS \textit{reg-to-reg} của \textbf{MAC-F} lên \textbf{$+3{,}524$\,ns}
($F_{\max,\mathrm{R2R}} = 154{,}42$\,MHz, tăng $+16{,}15$\,MHz) và WNS toàn thiết
kế đạt $+2{,}247$\,ns (tương ứng $F_{\max}$ toàn thiết kế $= 129{,}0$\,MHz trên
đường tổ hợp ra \texttt{in\_ready}). Trên cùng khối S1 dùng chung, \texttt{MAC-NF}
đạt WNS toàn thiết kế $+0{,}669$\,ns ($107{,}2$\,MHz) và WNS nội bộ $+2{,}639$\,ns
($135{,}85$\,MHz, giảm nhẹ $0{,}304$\,ns so với trước tách S1 do dao động đặt--định
tuyến ngẫu nhiên của một lần chạy P\&R đơn lẻ khi đường găng chuyển về nhánh tích
đến \texttt{u\_product\_pack} chứ không phải nghẽn tại S1; ngược lại, mức cải thiện
$+0{,}756$\,ns của \texttt{MAC-F} gắn trực tiếp với việc giải tỏa đường găng S1).
Như vậy, \textbf{MAC-F} vượt trội đồng thời ở tần số nội bộ ($+18{,}57$\,MHz),
tần số toàn thiết kế ($+21{,}8$\,MHz), độ trễ xử lý ($-3$ chu kỳ/MAC) và tài nguyên LUT/FF.

\section{Phân tích kết quả triển khai MAC-Quire}

Quá trình tối ưu \texttt{MAC-Quire} gồm hai pha rõ rệt. Thứ nhất, chuyển đổi kiến
trúc backend từ dịch toàn vectơ 972 bit sang cập nhật cục bộ theo từng từ 36 bit
giúp cắt giảm 71,6\% LUT (10.156 xuống 2.888), 40,9\% FF (4.458 xuống 2.636) và
gần một nửa độ trễ F2F sau tổng hợp. Thứ hai, khi định tuyến đầy đủ OOC 100\,MHz
(4.594 net sạch, 0 lỗi DRC nghiêm trọng), mẫu \texttt{MAC-Quire} đạt $+0{,}170$\,ns
sau tổng hợp nhưng chuyển sang vi phạm setup nhẹ $-0{,}187$\,ns sau định tuyến
($\mathrm{TNS}=-19{,}336$\,ns, $F_{\max,\mathrm{R2R}} = 98{,}16$\,MHz; hold đạt
$+0{,}040$\,ns). Đường găng dài $9{,}851$\,ns (17 tầng logic, gồm 9 tầng chuỗi cộng)
rơi vào mạch giải mã địa chỉ từ tới chân cho phép của ngân hàng thanh ghi quire
(Bảng~\ref{tab:route-critical}). Như vậy, khi cần triển khai IP chạy trên $100$\,MHz
với tài nguyên gọn (ít hơn 2,35 lần LUT và 3,38 lần FF), \textbf{MAC-F} là lựa chọn
tối ưu hơn, trong khi \texttt{MAC-Quire} phù hợp cho ứng dụng triệt tiêu dải rộng
chấp nhận xung nhịp thấp hơn hoặc chèn thêm tầng thanh ghi ở đường chọn từ.

